\begin{thebibliography}{999}

\bibitem{bakcoleman2025actionable}
Joseph B. Bak-Coleman, Stephan Lewandowsky, Philipp Lorenz-Spreen, Arvind Narayanan, Amy Orben, and Lisa Oswald. (2025). Moving towards informative and actionable social media research. \textit{arXiv preprint arXiv:2505.09254v3}. \url{https://doi.org/10.48550/arxiv.2505.09254}.

\bibitem{bakcoleman2025rct}
Joseph B. Bak-Coleman. (2025). Randomized controlled trials cannot identify social media's causal effects. \textit{Nature Reviews Psychology}, \textit{4}, 615. \url{https://doi.org/10.1038/s44159-025-00475-5}.

\bibitem{sznajdweron2024toward}
Katarzyna Sznajd-Weron, Arkadiusz J\k{e}drzejewski, and Barbara Kami\'{n}ska. (2024). Toward Understanding of the Social Hysteresis: Insights From Agent-Based Modeling. \textit{Perspectives on Psychological Science}, \textit{19}(2), 511-521. \url{https://doi.org/10.1177/17456916231195361}.

\bibitem{oestereich2020hysteresis}
Andr\'{e} L. Oestereich, Marcelo A. Pires, S\'{\i}lvio M. Duarte Queir\'{o}s, and Nuno Crokidakis. (2020). Hysteresis and disorder-induced order in continuous kinetic-like opinion dynamics in complex networks. \textit{Chaos Solitons \& Fractals}, \textit{137}, 109893. \url{https://doi.org/10.1016/j.chaos.2020.109893}.

\bibitem{nowak2020threshold}
B. Nowak, and Katarzyna Sznajd-Weron. (2020). Symmetrical threshold model with independence on random graphs. \textit{Physical review. E}, \textit{101}(5), 052316. \url{https://doi.org/10.1103/physreve.101.052316}.

\bibitem{macy2021polarization}
Michael W. Macy, Manqing Ma, Daniel R. Tabin, Jianxi Gao, and Boles\l{}aw K. Szyma\'{n}ski. (2021). Polarization and tipping points. \textit{Proceedings of the National Academy of Sciences}, \textit{118}(50), e2102144118. \url{https://doi.org/10.1073/pnas.2102144118}.

\bibitem{su2017contagion}
Zhen Su, Wei Wang, Lixiang Li, Jinghua Xiao, and H. Eugene Stanley. (2017). Emergence of hysteresis loop in social contagions on complex networks. \textit{Scientific Reports}, \textit{7}(1), 6103. \url{https://doi.org/10.1038/s41598-017-06286-w}.

\bibitem{chen2019vaccine}
Xingru Chen, and Feng Fu. (2019). Imperfect vaccine and hysteresis. \textit{Proceedings of the Royal Society B Biological Sciences}, \textit{286}(1894), 20182406. \url{https://doi.org/10.1098/rspb.2018.2406}.

\bibitem{graham2020sociology}
Hannah Graham. (2020). Hysteresis and the sociological perspective in a time of crisis. \textit{Acta Sociologica}, \textit{63}(4), 450-452. \url{https://doi.org/10.1177/0001699320961814}.

\bibitem{maia2021adaptive}
Hugo P. Maia, Silvio C. Ferreira, and Marcelo L. Martins. (2021). Adaptive network approach for emergence of societal bubbles. \textit{Physica A Statistical Mechanics and its Applications}, \textit{572}, 125588. \url{https://doi.org/10.1016/j.physa.2020.125588}.

\bibitem{asikainen2020cumulative}
Aili Asikainen, Gerardo I\~{n}iguez, Javier Ure\~{n}a-Carri\'{o}n, Kimmo Kaski, and Mikko Kivel\"{a}. (2020). Cumulative effects of triadic closure and homophily in social networks. \textit{Science Advances}, \textit{6}(19), eaax7310. \url{https://doi.org/10.1126/sciadv.aax7310}.

\bibitem{zeng2015mutual}
An Zeng, Chi Ho Yeung, Mat\'{u}\v{s} Medo, and Yi-Cheng Zhang. (2015). Modeling mutual feedback between users and recommender systems. \textit{Journal of Statistical Mechanics Theory and Experiment}, \textit{2015}(7), P07020. \url{https://doi.org/10.1088/1742-5468/2015/07/p07020}.

\bibitem{chan2026brake}
Johnny Chan. (2026). Where the Brake Slips: A Feedback View of the Micro-Drama Platform. \textit{Systems Research and Behavioral Science}, 1-7. \url{https://doi.org/10.1002/sres.70126}.

\bibitem{davidson2026closedloop}
Ella C. Davidson, and Mengbin Ye. (2026). Modelling the Closed-Loop Dynamics Between a Social Media Recommender System and Users' Opinions. \textit{IEEE Transactions on Computational Social Systems}, 1--13. \url{https://doi.org/10.1109/TCSS.2026.3715162}.

\bibitem{depasquale2026control}
Giulia De Pasquale, Sarah Dean, and Paolo Frasca. (2026). Recommender Systems as Control Systems. \textit{arXiv preprint arXiv:2605.01503}. \url{https://doi.org/10.48550/arxiv.2605.01503}.

\bibitem{mariano2026optimal}
Simone Mariano, and Paolo Frasca. (2026). Optimal Control Synthesis of Closed-Loop Recommendation Systems over Social Networks. \textit{arXiv preprint arXiv:2603.10275}. \url{https://doi.org/10.48550/arxiv.2603.10275}.

\bibitem{mariano2026pathological}
Simone Mariano, and Paolo Frasca. (2026). Pathological Regimes of Closed-Loop Recommendation Systems over Social Networks. \textit{arXiv preprint arXiv:2607.05010}. \url{https://doi.org/10.48550/arxiv.2607.05010}.

\bibitem{gauthier2026x}
Germain Gauthier, Roland Hodler, Philine Widmer, and Ekaterina Zhuravskaya. (2026). The political effects of X's feed algorithm. \textit{Nature}, \textit{652}(8109), 416-423. \url{https://doi.org/10.1038/s41586-026-10098-2}.

\bibitem{guess2023feeds}
Andrew M. Guess, Neil Malhotra, Jennifer Pan, Pablo Barber\'{a}, Hunt Allcott, Taylor Brown, Adriana Crespo-Tenorio, Drew Dimmery, Deen Freelon, Matthew Gentzkow, Sandra Gonz\'{a}lez-Bail\'{o}n, Edward H. Kennedy, Young Mie Kim, David Lazer, Devra Moehler, Brendan Nyhan, Carlos Velasco Rivera, Jaime E. Settle, Daniel Robert Thomas, ... , and Joshua A. Tucker. (2023). How do social media feed algorithms affect attitudes and behavior in an election campaign?. \textit{Science}, \textit{381}(6656), 398-404. \url{https://doi.org/10.1126/science.abp9364}.

\bibitem{nyhan2023likeminded}
Brendan Nyhan, Jaime E. Settle, Emily Thorson, Magdalena Wojcieszak, Pablo Barber\'{a}, Annie Y. Chen, Hunt Allcott, Taylor Brown, Adriana Crespo-Tenorio, Drew Dimmery, Deen Freelon, Matthew Gentzkow, Sandra Gonz\'{a}lez-Bail\'{o}n, Andrew M. Guess, Edward H. Kennedy, Young Mie Kim, David Lazer, Neil Malhotra, Devra Moehler, ... , and Joshua A. Tucker. (2023). Like-minded sources on Facebook are prevalent but not polarizing. \textit{Nature}, \textit{620}(7972), 137-144. \url{https://doi.org/10.1038/s41586-023-06297-w}.

\bibitem{gonzalez2023segregation}
Sandra Gonz\'{a}lez-Bail\'{o}n, David Lazer, Pablo Barber\'{a}, Meiqing Zhang, Hunt Allcott, Taylor Brown, Adriana Crespo-Tenorio, Deen Freelon, Matthew Gentzkow, Andrew M. Guess, Shanto Iyengar, Young Mie Kim, Neil Malhotra, Devra Moehler, Brendan Nyhan, Jennifer Pan, Carlos Velasco Rivera, Jaime E. Settle, Emily Thorson, ... , and Joshua A. Tucker. (2023). Asymmetric ideological segregation in exposure to political news on Facebook. \textit{Science}, \textit{381}(6656), 392-398. \url{https://doi.org/10.1126/science.ade7138}.

\bibitem{guess2023reshares}
Andrew M. Guess, Neil Malhotra, Jennifer Pan, Pablo Barber\'{a}, Hunt Allcott, Taylor Brown, Adriana Crespo-Tenorio, Drew Dimmery, Deen Freelon, Matthew Gentzkow, Sandra Gonz\'{a}lez-Bail\'{o}n, Edward H. Kennedy, Young Mie Kim, David Lazer, Devra Moehler, Brendan Nyhan, Carlos Velasco Rivera, Jaime E. Settle, Daniel Robert Thomas, ... , and Joshua A. Tucker. (2023). Reshares on social media amplify political news but do not detectably affect beliefs or opinions. \textit{Science}, \textit{381}(6656), 404-408. \url{https://doi.org/10.1126/science.add8424}.

\bibitem{centola2018tipping}
Damon Centola, Joshua Becker, Devon Brackbill, and Andrea Baronchelli. (2018). Experimental evidence for tipping points in social convention. \textit{Science}, \textit{360}(6393), 1116-1119. \url{https://doi.org/10.1126/science.aas8827}.

\bibitem{scheffer2009warning}
Marten Scheffer, Jordi Bascompte, William A. Brock, Victor Brovkin, Stephen R. Carpenter, Vasilis Dakos, Hermann Held, Egbert H. van Nes, Max Rietkerk, and George Sugihara. (2009). Early-warning signals for critical transitions. \textit{Nature}, \textit{461}(7260), 53-59. \url{https://doi.org/10.1038/nature08227}.

\bibitem{wiedermann2020threshold}
Marc Wiedermann, E. Keith Smith, Jobst Heitzig, and Jonathan F. Donges. (2020). A network-based microfoundation of Granovetter's threshold model for social tipping. \textit{Scientific Reports}, \textit{10}(1), 11202. \url{https://doi.org/10.1038/s41598-020-67102-6}.

\bibitem{kaminska2025cognitive}
Barbara Kami\'{n}ska, and Katarzyna Sznajd-Weron. (2025). Impact of cognitive dissonance on social hysteresis: Insights from the expressed and private opinions model. \textit{Expert Systems with Applications}, \textit{273}, 126851. \url{https://doi.org/10.1016/j.eswa.2025.126851}.

\bibitem{cruzcastillo2026spiral}
Juan Guillermo Cruz-Castillo, and Emanuele Cozzo. (2026). Modelling Public Opinion Dynamics: The Spiral of silence in clustered homophilic networks. \textit{Chaos Solitons \& Fractals}, \textit{203}, 117660. \url{https://doi.org/10.1016/j.chaos.2025.117660}.

\bibitem{gross2008adaptive}
Thilo Gross, and Bernd Blasius. (2007). Adaptive coevolutionary networks: a review. \textit{Journal of The Royal Society Interface}, \textit{5}(20), 259-271. \url{https://doi.org/10.1098/rsif.2007.1229}.

\bibitem{holme2006nonequilibrium}
Petter Holme, and M. E. J. Newman. (2006). Nonequilibrium phase transition in the coevolution of networks and opinions. \textit{Physical Review E}, \textit{74}(5), 056108. \url{https://doi.org/10.1103/physreve.74.056108}.

\bibitem{kozma2008consensus}
B. Kozma, and Alain Barrat. (2008). Consensus formation on adaptive networks. \textit{Physical Review E}, \textit{77}(1), 016102. \url{https://doi.org/10.1103/physreve.77.016102}.

\bibitem{gil2006coevolution}
Santiago Gil, and Dami\'{a}n H. Zanette. (2006). Coevolution of agents and networks: Opinion spreading and community disconnection. \textit{Physics Letters A}, \textit{356}(2), 89-94. \url{https://doi.org/10.1016/j.physleta.2006.03.037}.

\bibitem{vazquez2008adaptive}
Federico V\'{a}zquez, Victor M. Egu\'{\i}luz, and M. San Miguel. (2008). Generic Absorbing Transition in Coevolution Dynamics. \textit{Physical Review Letters}, \textit{100}(10), 108702. \url{https://doi.org/10.1103/physrevlett.100.108702}.

\bibitem{kan2023adaptive}
Unchitta Kan, Michelle Feng, and Mason A. Porter. (2022). An adaptive bounded-confidence model of opinion dynamics on networks. \textit{Journal of Complex Networks}, \textit{11}(1), cnac055. \url{https://doi.org/10.1093/comnet/cnac055}.

\bibitem{li2025adaptive}
Grace J. Li, Jiajie Luo, and Mason A. Porter. (2025). Bounded-Confidence Models of Opinion Dynamics with Adaptive Confidence Bounds. \textit{SIAM Journal on Applied Dynamical Systems}, \textit{24}(2), 994-1041. \url{https://doi.org/10.1137/23m1558951}.

\bibitem{deffuant2000mixing}
Guillaume Deffuant, David B. Neau, Fr\'{e}d\'{e}ric Amblard, and G\'{e}rard Weisbuch. (2000). Mixing beliefs among interacting agents. \textit{Advances in Complex Systems}, \textit{03}(01n04), 87-98. \url{https://doi.org/10.1142/s0219525900000078}.

\bibitem{castellano2009statistical}
Claudio Castellano, Santo Fortunato, and Vittorio Loreto. (2009). Statistical physics of social dynamics. \textit{Reviews of Modern Physics}, \textit{81}(2), 591-646. \url{https://doi.org/10.1103/revmodphys.81.591}.

\bibitem{dandekar2013assimilation}
Pranav Dandekar, Ashish Goel, and David T. Lee. (2013). Biased assimilation, homophily, and the dynamics of polarization. \textit{Proceedings of the National Academy of Sciences}, \textit{110}(15), 5791-5796. \url{https://doi.org/10.1073/pnas.1217220110}.

\bibitem{sasahara2021unfollowing}
Kazutoshi Sasahara, Wen Chen, Hao Peng, Giovanni Luca Ciampaglia, Alessandro Flammini, and Filippo Menczer. (2020). Social influence and unfollowing accelerate the emergence of echo chambers. \textit{Journal of Computational Social Science}, \textit{4}(1), 381-402. \url{https://doi.org/10.1007/s42001-020-00084-7}.

\bibitem{wang2020public}
Xin Wang, Antonio Sirianni, Shaoting Tang, Zhiming Zheng, and Feng Fu. (2020). Public Discourse and Social Network Echo Chambers Driven by Socio-Cognitive Biases. \textit{Physical Review X}, \textit{10}(4), 041042. \url{https://doi.org/10.1103/physrevx.10.041042}.

\bibitem{cinelli2021echo}
Matteo Cinelli, Gianmarco De Francisci Morales, Alessandro Galeazzi, Walter Quattrociocchi, and Michele Starnini. (2021). The echo chamber effect on social media. \textit{Proceedings of the National Academy of Sciences}, \textit{118}(9), e2023301118. \url{https://doi.org/10.1073/pnas.2023301118}.

\bibitem{baumann2020echo}
Fabian Baumann, Philipp Lorenz-Spreen, Igor M. Sokolov, and Michele Starnini. (2020). Modeling Echo Chambers and Polarization Dynamics in Social Networks. \textit{Physical Review Letters}, \textit{124}(4), 048301. \url{https://doi.org/10.1103/physrevlett.124.048301}.

\bibitem{baumann2021topics}
Fabian Baumann, Philipp Lorenz-Spreen, Igor M. Sokolov, and Michele Starnini. (2021). Emergence of Polarized Ideological Opinions in Multidimensional Topic Spaces. \textit{Physical Review X}, \textit{11}(1), 011012. \url{https://doi.org/10.1103/physrevx.11.011012}.

\bibitem{thurner2025interactions}
Stefan Thurner, Markus Hofer, and Jan Korbel. (2025). Why more social interactions lead to more polarization in societies. \textit{Proceedings of the National Academy of Sciences}, \textit{122}(44), e2517530122. \url{https://doi.org/10.1073/pnas.2517530122}.

\bibitem{preisach1935magnetic}
F. Preisach. (1935). \"{U}ber die magnetische Nachwirkung. \textit{Zeitschrift f\"{u}r Physik}, \textit{94}(5-6), 277-302. \url{https://doi.org/10.1007/bf01349418}.

\bibitem{macki1993mathematical}
Jack W. Macki, Paolo Nistri, and Pietro Zecca. (1993). Mathematical Models for Hysteresis. \textit{SIAM Review}, \textit{35}(1), 94-123. \url{https://doi.org/10.1137/1035005}.

\bibitem{hassani2014survey}
Vahid Hassani, Tegoeh Tjahjowidodo, and Thanh Nho. (2014). A survey on hysteresis modeling, identification and control. \textit{Mechanical Systems and Signal Processing}, \textit{49}(1-2), 209-233. \url{https://doi.org/10.1016/j.ymssp.2014.04.012}.

\bibitem{oh2005duhem}
JinHyoung Oh, and Dennis S. Bernstein. (2005). Semilinear Duhem model for rate-independent and rate-dependent hysteresis. \textit{IEEE Transactions on Automatic Control}, \textit{50}(5), 631-645. \url{https://doi.org/10.1109/tac.2005.847035}.

\bibitem{sethna1993hierarchies}
James P. Sethna, Karin A. Dahmen, Sivan Kartha, J. A. Krumhansl, Bruce W. Roberts, and Joel D. Shore. (1993). Hysteresis and hierarchies: Dynamics of disorder-driven first-order phase transformations. \textit{Physical Review Letters}, \textit{70}(21), 3347-3350. \url{https://doi.org/10.1103/physrevlett.70.3347}.

\bibitem{mungan2019returnpoint}
Muhittin Mungan, and M. Mert Terzi. (2019). The Structure of State Transition Graphs in Systems with Return Point Memory: I. General Theory. \textit{Annales Henri Poincar\'{e}}, \textit{20}(8), 2819-2872. \url{https://doi.org/10.1007/s00023-019-00807-1}.

\bibitem{krejci2001inverse}
Pavel Krej\v{c}\'{\i}, and Klaus Kuhnen. (2001). Inverse control of systems with hysteresis and creep. \textit{IEE Proceedings---Control Theory and Applications}, \textit{148}(3), 185-192. \url{https://doi.org/10.1049/ip-cta:20010375}.

\bibitem{ikhouane2007systems}
Fay\c{c}al Ikhouane, and Jos\'{e} Rodellar. (2007). Systems with Hysteresis. \textit{John Wiley \& Sons}. \url{https://doi.org/10.1002/9780470513200}.

\bibitem{perra2019personalisation}
Nicola Perra, and Luis E. C. Rocha. (2019). Modelling opinion dynamics in the age of algorithmic personalisation. \textit{Scientific Reports}, \textit{9}(1), 7261. \url{https://doi.org/10.1038/s41598-019-43830-2}.

\bibitem{sirbu2019bias}
Alina S\^{i}rbu, Dino Pedreschi, Fosca Giannotti, and J\'{a}nos Kert\'{e}sz. (2019). Algorithmic bias amplifies opinion fragmentation and polarization: A bounded confidence model. \textit{PLoS ONE}, \textit{14}(3), e0213246. \url{https://doi.org/10.1371/journal.pone.0213246}.

\bibitem{arruda2022algorithms}
Henrique Ferraz de Arruda, Felipe Maciel Cardoso, Guilherme Ferraz de Arruda, Alexis Hern\'{a}ndez, Luciano da Fontoura Costa, and Yamir Moreno. (2022). Modelling how social network algorithms can influence opinion polarization. \textit{Information Sciences}, \textit{588}, 265-278. \url{https://doi.org/10.1016/j.ins.2021.12.069}.

\bibitem{peralta2020asymmetry}
Giordano De Marzo, Andrea Zaccaria, and Claudio Castellano. (2020). Emergence of polarization in a voter model with personalized information. \textit{Physical Review Research}, \textit{2}(4), 043117. \url{https://doi.org/10.1103/physrevresearch.2.043117}.

\bibitem{valensise2023drivers}
Carlo Michele Valensise, Matteo Cinelli, and Walter Quattrociocchi. (2023). The drivers of online polarization: Fitting models to data. \textit{Information Sciences}, \textit{642}, 119152. \url{https://doi.org/10.1016/j.ins.2023.119152}.

\bibitem{donkers2024desounding}
Tim Donkers, and J\"{u}rgen Ziegler. (2023). De-sounding echo chambers: Simulation-based analysis of polarization dynamics in social networks. \textit{Online Social Networks and Media}, \textit{37-38}, 100275. \url{https://doi.org/10.1016/j.osnem.2023.100275}.

\bibitem{rossi2018closedloop}
Wilbert Samuel Rossi, Jan Willem Polderman, and Paolo Frasca. (2018). The closed loop between opinion formation and personalised recommendations. \textit{arXiv preprint arXiv:1809.04644}. \url{https://doi.org/10.48550/arxiv.1809.04644}.

\bibitem{gezha2026control}
Vladislav N. Gezha, and Ivan Kozitsin. (2026). Control over Recommendation Algorithms in Heterogeneous Modular Systems with Dynamic Opinions. \textit{Entropy}, \textit{28}(3), 333. \url{https://doi.org/10.3390/e28030333}.

\bibitem{oliveira2026interplay}
Kleber A. Oliveira, Henrique Ferraz de Arruda, and Yamir Moreno. (2026). Mechanistic interplay between information spreading and opinion polarization. \textit{PNAS Nexus}, \textit{5}(1), pgaf402. \url{https://doi.org/10.1093/pnasnexus/pgaf402}.

\bibitem{cinus2022people}
Federico Cinus, Marco Minici, Corrado Monti, and Francesco Bonchi. (2022). The Effect of People Recommenders on Echo Chambers and Polarization. \textit{Proceedings of the International AAAI Conference on Web and Social Media}, \textit{16}, 90-101. \url{https://doi.org/10.1609/icwsm.v16i1.19275}.

\bibitem{chaney2018confounding}
Allison J. B. Chaney, Brandon Stewart, and Barbara E. Engelhardt. (2018). How algorithmic confounding in recommendation systems increases homogeneity and decreases utility. \textit{Proceedings of the 12th ACM Conference on Recommender Systems}, 224-232. \url{https://doi.org/10.1145/3240323.3240370}.

\bibitem{mansoury2020feedback}
Masoud Mansoury, Himan Abdollahpouri, Mykola Pechenizkiy, Bamshad Mobasher, and Robin Burke. (2020). Feedback Loop and Bias Amplification in Recommender Systems. \textit{Proceedings of the 29th ACM International Conference on Information and Knowledge Management}, 2145-2148. \url{https://doi.org/10.1145/3340531.3412152}.

\bibitem{kalimeris2021amplification}
Dimitris Kalimeris, Smriti Bhagat, Shankar Kalyanaraman, and Udi Weinsberg. (2021). Preference Amplification in Recommender Systems. \textit{Proceedings of the 27th ACM SIGKDD Conference on Knowledge Discovery and Data Mining}, 805-815. \url{https://doi.org/10.1145/3447548.3467298}.

\bibitem{zhao2013longterm}
Dan-Dan Zhao, An Zeng, Mingsheng Shang, and Jian Gao. (2013). Long-Term Effects of Recommendation on the Evolution of Online Systems. \textit{Chinese Physics Letters}, \textit{30}(11), 118901. \url{https://doi.org/10.1088/0256-307x/30/11/118901}.

\bibitem{su2016network}
Jessica Su, Aneesh Sharma, and Sharad Goel. (2016). The Effect of Recommendations on Network Structure. \textit{Proceedings of the 25th International Conference on World Wide Web}, 1157-1167. \url{https://doi.org/10.1145/2872427.2883040}.

\bibitem{coppolillo2025drift}
Erica Coppolillo, Simone Mungari, Ettore Ritacco, Francesco Fabbri, Marco Minici, Francesco Bonchi, and Giuseppe Manco. (2025). Algorithmic Drift: A simulation framework to study the effects of recommender systems on user preferences. \textit{Information Processing \& Management}, \textit{62}(4), 104125. \url{https://doi.org/10.1016/j.ipm.2025.104125}.

\bibitem{chandrasekaran2026networkaware}
Sanjay Chandrasekaran, Giulia De Pasquale, Giuseppe Belgioioso, and Florian D\"{o}rfler. (2026). Network-Aware Recommender System via Online Feedback Optimization. \textit{IEEE Transactions on Automatic Control}, \textit{71}(3), 1708-1723. \url{https://doi.org/10.1109/tac.2025.3616262}.

\bibitem{moller2018diversity}
Judith M\"{o}ller, Damian Trilling, Natali Helberger, and Bram van Es. (2018). Do not blame it on the algorithm: an empirical assessment of multiple recommender systems and their impact on content diversity. \textit{Information Communication \& Society}, \textit{21}(7), 959-977. \url{https://doi.org/10.1080/1369118x.2018.1444076}.

\bibitem{fabbri2022minorities}
Antonio Ferrara, Lisette Esp\'{i}n-Noboa, Fariba Karimi, and Claudia Wagner. (2022). Link recommendations: Their impact on network structure and minorities. \textit{Proceedings of the 14th ACM Web Science Conference}, 228-238. \url{https://doi.org/10.1145/3501247.3531583}.

\bibitem{frey2010connections}
Elizabeth Daly, Werner Geyer, and David R. Millen. (2010). The network effects of recommending social connections. \textit{Proceedings of the Fourth ACM Conference on Recommender Systems}, 301-304. \url{https://doi.org/10.1145/1864708.1864772}.

\bibitem{jadidinejad2020exploration}
Amir H. Jadidinejad, Craig Macdonald, and Iadh Ounis. (2020). Using Exploration to Alleviate Closed Loop Effects in Recommender Systems. \textit{Proceedings of the 43rd International ACM SIGIR Conference on Research and Development in Information Retrieval}, 2025-2028. \url{https://doi.org/10.1145/3397271.3401230}.

\bibitem{krauth2025causal}
Karl Krauth, Yixin Wang, and Michael I. Jordan. (2025). Breaking Feedback Loops in Recommender Systems with Causal Inference. \textit{ACM Transactions on Recommender Systems}, \textit{4}(1), 1-20. \url{https://doi.org/10.1145/3728372}.

\bibitem{sinha2016deconvolving}
Ayan Sinha, David F. Gleich, and Karthik Ramani. (2017). Deconvolving Feedback Loops in Recommender Systems. \textit{arXiv preprint arXiv:1703.01049}. \url{https://doi.org/10.48550/arxiv.1703.01049}.

\bibitem{jannach2020mcnamara}
Dietmar Jannach, and Christine Bauer. (2020). Escaping the McNamara Fallacy: Toward More Impactful Recommender Systems Research. \textit{AI Magazine}, \textit{41}(4), 79-95. \url{https://doi.org/10.1609/aimag.v41i4.5312}.

\bibitem{ge2020ecommerce}
Yingqiang Ge, Shuya Zhao, Honglu Zhou, Changhua Pei, Fei Sun, Wenwu Ou, and Yongfeng Zhang. (2020). Understanding Echo Chambers in E-commerce Recommender Systems. \textit{Proceedings of the 43rd International ACM SIGIR Conference on Research and Development in Information Retrieval}, 2261-2270. \url{https://doi.org/10.1145/3397271.3401431}.

\bibitem{zoralioglu2026dynamic}
Yildiz Zoralioglu, and Emre Yalcin. (2026). Dynamic feedback loops in recommender systems: Analyzing fairness, popularity bias, and user group disparities. \textit{Journal of Intelligent Information Systems}. \url{https://doi.org/10.1007/s10844-026-01025-y}.

\bibitem{huszar2022amplification}
Ferenc Husz\'{a}r, Sofia Ira Ktena, Conor Cruise O'Brien, Luca Belli, Andrew Schlaikjer, and Moritz Hardt. (2022). Algorithmic Amplification of Politics on Twitter. \textit{Proceedings of the National Academy of Sciences}, \textit{119}(1), e2025334119. \url{https://doi.org/10.1073/pnas.2025334119}.

\bibitem{bakshy2015exposure}
Eytan Bakshy, Solomon Messing, and Lada A. Adamic. (2015). Exposure to ideologically diverse news and opinion on Facebook. \textit{Science}, \textit{348}(6239), 1130-1132. \url{https://doi.org/10.1126/science.aaa1160}.

\bibitem{flaxman2016bubbles}
Seth Flaxman, Sharad Goel, and Justin M. Rao. (2016). Filter Bubbles, Echo Chambers, and Online News Consumption. \textit{Public Opinion Quarterly}, \textit{80}(S1), 298-320. \url{https://doi.org/10.1093/poq/nfw006}.

\bibitem{bandy2021twitter}
Jack Bandy, and Nicholas Diakopoulos. (2021). Curating Quality? How Twitter's Timeline Algorithm Treats Different Types of News. \textit{Social Media + Society}, \textit{7}(3), 20563051211041648. \url{https://doi.org/10.1177/20563051211041648}.

\bibitem{ibrahim2026tiktok}
Hazem Ibrahim, HyunSeok Daniel Jang, Nouar AlDahoul, Aaron R. Kaufman, Talal Rahwan, and Yasir Zaki. (2026). Systematic partisan content skews in TikTok during the 2024 US elections. \textit{Nature}, \textit{654}(8120), 1004-1011. \url{https://doi.org/10.1038/s41586-026-10447-1}.

\bibitem{hosseinmardi2024youtube}
Homa Hosseinmardi, Amir Ghasemian, Miguel Rivera-Lanas, Manoel Horta Ribeiro, Robert West, and Duncan J. Watts. (2024). Causally estimating the effect of YouTube's recommender system using counterfactual bots. \textit{Proceedings of the National Academy of Sciences}, \textit{121}(8), e2313377121. \url{https://doi.org/10.1073/pnas.2313377121}.

\bibitem{cho2020search}
Jaeho Cho, Saifuddin Ahmed, Martin Hilbert, Billy Liu, and Jonathan Luu. (2020). Do Search Algorithms Endanger Democracy? An Experimental Investigation of Algorithm Effects on Political Polarization. \textit{Journal of Broadcasting \& Electronic Media}, \textit{64}(2), 150-172. \url{https://doi.org/10.1080/08838151.2020.1757365}.

\bibitem{neumann2023curation}
Ole Kelm, Tim Neumann, Maike Behrendt, Markus Brenneis, Katharina Gerl, Stefan Marschall, Florian Mei{\ss}ner, Stefan Harmeling, Gerhard Vowe, and Marc Ziegele. (2023). How algorithmically curated online environments influence users' political polarization: Results from two experiments with panel data. \textit{Computers in Human Behavior Reports}, \textit{12}, 100343. \url{https://doi.org/10.1016/j.chbr.2023.100343}.

\bibitem{bail2018opposing}
Christopher A. Bail, Lisa P. Argyle, Taylor Brown, John P. Bumpus, Haohan Chen, M. B. Fallin Hunzaker, Jaemin Lee, Marcus Mann, Friedolin Merhout, and Alexander Volfovsky. (2018). Exposure to opposing views on social media can increase political polarization. \textit{Proceedings of the National Academy of Sciences}, \textit{115}(37), 9216-9221. \url{https://doi.org/10.1073/pnas.1804840115}.

\bibitem{piccardi2025reranking}
Tiziano Piccardi, Martin Saveski, Chenyan Jia, Jeffrey T. Hancock, Jeanne L. Tsai, and Michael S. Bernstein. (2025). Reranking partisan animosity in algorithmic social media feeds alters affective polarization. \textit{Science}, \textit{390}(6776), eadu5584. \url{https://doi.org/10.1126/science.adu5584}.

\bibitem{brady2026norms}
William J. Brady, Meriel Doyle, Abdo Elnakouri, Eli J. Finkel, Joshua Conrad Jackson, Nour Kteily, Victoria Parker, Curtis Puryear, Trevor Spelman, Jacob Teeny, and Mark Torres. (2026). Redesigning algorithms to intervene on social norm misperceptions during a national election. \textit{Nature}, \textit{655}(8124), 942-956. \url{https://doi.org/10.1038/s41586-026-10536-1}.

\bibitem{muchnik2013influence}
Lev Muchnik, Sinan Aral, and Sean J. Taylor. (2013). Social Influence Bias: A Randomized Experiment. \textit{Science}, \textit{341}(6146), 647-651. \url{https://doi.org/10.1126/science.1240466}.

\bibitem{bond2012mobilization}
Robert M. Bond, Christopher J. Fariss, Jason Jones, Adam Kramer, Cameron Marlow, Jaime E. Settle, and James H. Fowler. (2012). A 61-million-person experiment in social influence and political mobilization. \textit{Nature}, \textit{489}(7415), 295-298. \url{https://doi.org/10.1038/nature11421}.

\bibitem{allcott2024deactivation}
Hunt Allcott, Matthew Gentzkow, Winter Mason, Arjun Wilkins, Pablo Barber\'{a}, Taylor Brown, Juan Carlos Cisneros, Adriana Crespo-Tenorio, Drew Dimmery, Deen Freelon, Sandra Gonz\'{a}lez-Bail\'{o}n, Andrew M. Guess, Young Mie Kim, David Lazer, Neil Malhotra, Devra Moehler, Sameer Nair-Desai, Houda Nait El Barj, Brendan Nyhan, ... , and Joshua A. Tucker. (2024). The effects of Facebook and Instagram on the 2020 election: A deactivation experiment. \textit{Proceedings of the National Academy of Sciences}, \textit{121}(21), e2321584121. \url{https://doi.org/10.1073/pnas.2321584121}.

\bibitem{kubin2021review}
Emily Kubin, and Christian von Sikorski. (2021). The role of (social) media in political polarization: a systematic review. \textit{Annals of the International Communication Association}, \textit{45}(3), 188-206. \url{https://doi.org/10.1080/23808985.2021.1976070}.

\bibitem{vanbavel2021shapes}
Jay Joseph Van Bavel, Steve Rathje, Elizabeth Harris, Claire Robertson, and Anni Sternisko. (2021). How social media shapes polarization. \textit{Trends in Cognitive Sciences}, \textit{25}(11), 913-916. \url{https://doi.org/10.1016/j.tics.2021.07.013}.

\bibitem{hartmann2025review}
David Hartmann, Sonja Mei Wang, Lena Pohlmann, and Bettina Berendt. (2025). A systematic review of echo chamber research: comparative analysis of conceptualizations, operationalizations, and varying outcomes. \textit{Journal of Computational Social Science}, \textit{8}(2), 52. \url{https://doi.org/10.1007/s42001-025-00381-z}.

\bibitem{mahmoudi2024review}
Amin Mahmoudi, Dariusz Jemielniak, and Leon Ciechanowski. (2024). Echo Chambers in Online Social Networks: A Systematic Literature Review. \textit{IEEE Access}, \textit{12}, 9594-9620. \url{https://doi.org/10.1109/access.2024.3353054}.

\bibitem{gao2024causal}
Huishi Luo, Fuzhen Zhuang, Ruobing Xie, Hengshu Zhu, Deqing Wang, Zhulin An, and Yongjun Xu. (2024). A survey on causal inference for recommendation. \textit{The Innovation}, \textit{5}(2), 100590. \url{https://doi.org/10.1016/j.xinn.2024.100590}.

\bibitem{ge2024trustworthy}
Yingqiang Ge, Shuchang Liu, Zuohui Fu, Juntao Tan, Zelong Li, Shuyuan Xu, Yunqi Li, Yikun Xian, and Yongfeng Zhang. (2024). A Survey on Trustworthy Recommender Systems. \textit{ACM Transactions on Recommender Systems}, \textit{3}(2), 1-68. \url{https://doi.org/10.1145/3652891}.

\bibitem{aridor2024economics}
Guy Aridor, Rafael Jim\'{e}nez-Dur\'{a}n, Roee Levy, and Lena Song. (2024). The Economics of Social Media. \textit{Journal of Economic Literature}, \textit{62}(4), 1422-1474. \url{https://doi.org/10.1257/jel.20241743}.

\bibitem{noorazar2020survey}
Hossein Noorazar. (2020). Recent advances in opinion propagation dynamics: a 2020 survey. \textit{The European Physical Journal Plus}, \textit{135}(6), 521. \url{https://doi.org/10.1140/epjp/s13360-020-00541-2}.

\bibitem{hudgens2008interference}
Michael G. Hudgens, and M. Elizabeth Halloran. (2008). Toward Causal Inference With Interference. \textit{Journal of the American Statistical Association}, \textit{103}(482), 832-842. \url{https://doi.org/10.1198/016214508000000292}.

\bibitem{aronow2017interference}
Peter M. Aronow, and Cyrus Samii. (2017). Estimating average causal effects under general interference, with application to a social network experiment. \textit{The Annals of Applied Statistics}, \textit{11}(4). \url{https://doi.org/10.1214/16-aoas1005}.

\bibitem{liu2026saturation}
Jizhou Liu, Azeem M. Shaikh, and Liang Zhong. (2026). Randomization Tests in Randomized Saturation Designs. \textit{arXiv preprint arXiv:2607.04257}. \url{https://doi.org/10.48550/arxiv.2607.04257}.

\bibitem{bojinov2023switchback}
Iavor Bojinov, David Simchi-Levi, and Jinglong Zhao. (2023). Design and Analysis of Switchback Experiments. \textit{Management Science}, \textit{69}(7), 3759-3777. \url{https://doi.org/10.1287/mnsc.2022.4583}.

\bibitem{wang2024carryover}
Danni Shi, and Ting Ye. (2024). Behavioral carry-over effect and power consideration in crossover trials. \textit{Biometrics}, \textit{80}(2), ujae023. \url{https://doi.org/10.1093/biomtc/ujae023}.

\bibitem{sun2025dynamic}
Jianping Sun, Peiran Guo, Xiaoyang Chen, and Xianming Tan. (2025). Dynamic Treatment Effect Analysis in Crossover Designs Through Repeated Measures. \textit{Statistics in Medicine}, \textit{44}(7), e70070. \url{https://doi.org/10.1002/sim.70070}.

\bibitem{wu2021randomization}
Jason Wu, and Peng Ding. (2021). Randomization Tests for Weak Null Hypotheses in Randomized Experiments. \textit{Journal of the American Statistical Association}, \textit{116}(536), 1898-1913. \url{https://doi.org/10.1080/01621459.2020.1750415}.

\bibitem{athey2019forests}
Susan Athey, Julie Tibshirani, and Stefan Wager. (2019). Generalized random forests. \textit{The Annals of Statistics}, \textit{47}(2). \url{https://doi.org/10.1214/18-aos1709}.

\end{thebibliography}
